% ============================================================
% 9. Exception Handling & Persistence
% ============================================================
\chapter{Exception Handling and Persistence}\label{ch:exceptions}

\section{Exception Handling}

Exceptions are modelled as a \emph{domain-specific} hierarchy that never
exposes physical persistence failures to the view layer:

\begin{itemize}
    \item \texttt{DAOException} (runtime): wraps physical failures such as
    \texttt{SQLException} raised by the JDBC layer, converting them into a
    domain-abstraction error at the DAO boundary;
    \item \texttt{IllegalArgumentException} / \texttt{IllegalStateException}:
    validation and business-rule violations raised by entities and controllers
    (e.g.\ invalid quantity, non-approved vendor, invalid order transition).
\end{itemize}

This design guarantees that the application flow is never interrupted by
leaked technical exceptions, and that every failure is reported through a
controlled, domain-meaningful message.

\section{Persistence}

The system supports \textbf{three persistence modes} (NFR-01), selected at
startup through the \texttt{ApplicationModeManager} and the
\texttt{DAOFactory} (Abstract Factory):

\begin{table}[H]
\centering
\caption{Persistence modes.}
\label{tab:persistence}
\begin{tabular}{@{}lll@{}}
\toprule
\textbf{Mode} & \textbf{Technology} & \textbf{Use} \\
\midrule
JDBC & MySQL (prepared statements, NFR-02) & production \\
File system & JSON via Gson (local-date adapters) & portable persistence \\
Demo & in-memory maps & tests and demo \\
\bottomrule
\end{tabular}
\end{table}

Every DAO (utente, prodotto, ricetta, ordine) is defined by an interface and
implemented three times; the factory guarantees that the whole family is
coherent. Password storage uses a one-way SHA-256 hash with a fixed salt
(NFR-03). The physical schema of the database is intentionally minimal: the
evaluation focuses on the interaction between controllers, entities and the
DAO layer.
